% ---
% id: reductions/req:emission-refusal
% kind: req
% title: "Refusal of an emitted family that is not exactly admissible"
% summary: >-
%   The engine emits a family only after deciding the six emission conditions
%   on the numbers about to be emitted, and otherwise refuses, repairing
%   nothing it was handed. A moved class, the flagged last resort included, is
%   refused only under the strict option, where it would rest on the snap.
% prominence: standard
% status: implemented
% uses: [def:emitnorm, thm:emit]
% mentions: [def:exactadm, req:reference-companion, foundations/req:snap-reference, foundations/req:support-safe-snap, req:emission-record, receiver/detector-imperfections/prop:noise-floor]
% assumes: []
% refs: []
% provenance: newly-derived
% implements: certiqs_key.reductions.validation:rule_emitted_family
% lean:
% ---

The decision on the six conditions is taken on what comes out of the
normalizer (Definition~\ref{def:emitnorm}\,(4), Theorem~\ref{thm:emit}\,(ii)).
A class the normalizer moves is not refused by default: it is snapped without
reviving a zero row or column and disclosed
(Requirement~\ref{req:emission-record}), as
\xref{foundations/req:snap-reference}\,(iii) asks of any snap of a
measurement family. A class whose residue no single element can take
without reviving a zero row or column has no support-safe snap in the
normalizer; it is snapped by the last resort, which always yields a positive
semidefinite family and is flagged in the record
(\xref{foundations/req:support-safe-snap}\,(iv)), so it is not refused by
default either. The usual cause of a last resort is a class whose elements
are rank-deficient on the sector, which a model with the device's loss and
dark counts can remove where the conditions the receiver module states for it
hold (\xref{receiver/detector-imperfections/prop:noise-floor}). One refusal
remains for a moved class: a run that declares the strict option refuses
wherever the class would rest on the snapped-measurement premise, which is
certain at emission where no reference companion that bypasses the
normalizer is formed (Requirement~\ref{req:reference-companion}).

\begin{requirement}\label{req:emission-refusal}
The engine emits a family only if the six conditions of
Theorem~\ref{thm:emit}\,(ii) --- (a)--(e) of
Definition~\ref{def:exactadm} together with \eqref{eq:latinv} of every entry
--- hold of the numbers about to be emitted, decided exactly after the
\tref{emission normalizer} where that acted; otherwise it refuses with the
reason \iface{emitted_family_not_admissible}, naming the predicate that
failed, and repairs nothing it was handed. Under the same reason, where the
run declares the strict option \iface{strict_snaps}
(\xref{foundations/req:snap-reference}\,(iv)) and no reference companion of
Requirement~\ref{req:reference-companion} is formed beside the artefact, it
refuses every moved class of Definition~\ref{def:emitnorm}, one snapped by
the last resort (3) included, naming the \tref{data sector}, the class, the
cause and the fix; where a companion is formed, whether the class rests is
decided by the run after its search, and the run refuses there under the
same option. Where the run does not declare the option, no class is refused
for being moved. The cause of a moved class that would rest is that no
reference companion bypasses the normalizer, which, by
Requirement~\ref{req:reference-companion}, is where no claim of this module
establishes (H-REF) for the instance; for a class snapped by (3) it also
names that a component of its residue has no host. The fix named is
declared numbers on which the exact construction stays on the lattice, so
that the class is kept; for a class
snapped by (3) also a model that includes the device's loss and dark counts,
so that the elements of the class are positive definite on the sector. A
moved class is mixed only with the family of its diagonal parts in (2), or
with a multiple of the identity in the last resort (3).
\end{requirement}
